\ifdefined\XeTeXrevision\else\pdfoutput=1\fi
\documentclass[11pt]{article}
\usepackage[margin=1in]{geometry}
\usepackage{booktabs}
\usepackage{amsmath}
\usepackage{xcolor}
\usepackage{hyperref}
\usepackage{microtype}
\hypersetup{
  colorlinks=true,
  linkcolor=blue!60!black,
  citecolor=blue!60!black,
  urlcolor=blue!60!black,
  pdftitle={Conditional LoRA Bridges at 230M: What Replicates, What Softens, and What Inverts},
  pdfauthor={Franci Penov},
}
\title{Conditional LoRA Bridges at 230M:\\What Replicates, What Softens, and What Inverts}
\author{Franci Penov \\ kortexa.ai \\ \texttt{francip@kortexa.ai}}
\date{2026-07-03}

\begin{document}
\maketitle

\begin{abstract}
Our controlled study of conditional LoRA bridges \cite{penov2026bridges} established its findings on a purpose-trained 33.6M-parameter model, leaving open which conclusions are properties of the mechanism and which are artifacts of scale. We port the unchanged experiment suite to a pretrained 230M-parameter model (LiquidAI LFM~2.5) and rerun it at both the standard and the $10\times$ training budget. The methodological core replicates exactly: bits-per-byte cannot detect sensor conditioning (the capacity-matched constant-feature control matches the true-feature bridge, +0.0852 vs +0.0856), while task-aligned probes detect it decisively (0.72 rank-1 on six-way IMU, 0.31 on fifty-way audio, all controls at chance). Three scale-dependent results reverse or soften. Prefix tuning flips sign, from -0.0048 BPB on the small model to +0.0614 at 230M, while remaining the weakest adapter. The diversity--BPB tradeoff softens: the strongest penalty retains 89\% of the unregularized BPB gain where the small model retained none. Most strikingly, the long-budget collapse inverts: extended unregularized task training destroyed IMU conditioning on the small model (rank-1 0.80 to 0.17, chance) but strengthens it at 230M (0.72 to 0.89) --- even as weight collapse under pure text loss persists at both scales. At the two scales tested, capacity and pretraining are allies of task-conditioned weight-space bridges, and single-scale claims about the mechanism's failure modes should be read as lower bounds.
\end{abstract}

\section{Motivation}
The conditional-LoRA-bridges study \cite{penov2026bridges} reported one sharp methodological result --- aggregate language-modeling metrics cannot certify sensor conditioning --- and one mechanism result, hypernetwork weight collapse, with facets that include degradation from token-space (prefix) conditioning and a budget-dependent destruction of task-level conditioning under prolonged unregularized training. All of it was measured on a single 33.6M-parameter model trained from scratch. This note answers the obvious question: which of those findings survive contact with a real pretrained model?

\section{Porting the harness}
We add a checkpoint adapter (\texttt{hf:<model-id>}) to the published harness so that every experiment --- baselines, controls, probes, budgets, seeds --- runs unchanged against a HuggingFace causal LM; here LiquidAI LFM~2.5 230M \cite{liquidai2025lfm2}, a 14-layer hybrid short-convolution/attention model with a 65{,}536-token vocabulary. Three porting details matter for fidelity:
\begin{itemize}
  \item \textbf{LoRA targets.} The paper's target set ``all'' (every attention and MLP linear) maps to q/k/v/out projections on the six attention layers and the SwiGLU projections on all fourteen layers; short-convolution projections are left unadapted. The generated LoRA vector has $D = 774{,}144$ dimensions at rank~4.
  \item \textbf{Byte-exact BPB across tokenizers.} Bits-per-byte is tokenizer-independent only if the per-token byte table reconstructs the exact UTF-8 length of any encoded text. We build and validate such a table for the model's byte-level BPE, including its non-special added tokens, so BPB values are comparable across the two vocabularies.
  \item \textbf{Sequence-start sensitivity.} LFM~2.5 is strongly BOS-dependent (7.13 nats/token without BOS vs.\ 2.70 with, on the same text). The packing dataloader already places BOS at every row start; the task probes prepend BOS for \texttt{hf:} models so the probe measures conditioning, not an out-of-distribution artifact. The published small-model protocol is unchanged and reproduces bit-identically.
\end{itemize}
The frozen 230M model evaluates at 1.1839 BPB on the paper's validation stream (the small model: 1.9654).

\section{What replicates}
Table~\ref{tab:bpb} gives the BPB decomposition at the standard budget. The paper's central methodological claim transfers unchanged: the conditional bridge (+0.0856) is indistinguishable from the capacity-matched constant-feature control (+0.0852), static LoRA is the best pure-text adapter (+0.0876), and three-seed spreads are below $10^{-3}$. BPB still cannot see conditioning; the task probes still can (Table~\ref{tab:task}): true features reach 0.72 rank-1 on six-way IMU and 0.31 on fifty-way audio while shuffled, random, and no-bridge controls sit at chance --- numerically close to the small model's 0.80 and 0.33. Composition of independently trained bridges remains sub-additive at both scales, and unregularized text-loss training still collapses generated weights toward input-independence (cross-input cosine 0.99 at $\lambda=0$).

\begin{table}[t]
\centering
\caption{BPB improvement over the frozen baseline (mean across audio, vision, IMU; standard budget; seed-42 protocol as in the original paper). The decomposition --- bridge $\approx$ constant control $<$ static LoRA --- replicates; prefix tuning flips sign.}
\label{tab:bpb}
\begin{tabular}{lrr}
\toprule
Method & 33.6M (paper) & 230M (this note) \\
\midrule
Static LoRA & +0.0033 & +0.0876 \\
Conditional bridge & +0.0027 & +0.0856 \\
Shuffled features & +0.0018 & +0.0853 \\
Random features & +0.0006 & +0.0823 \\
Constant features & +0.0027 & +0.0852 \\
Prefix tuning & -0.0048 & +0.0614 \\
\bottomrule
\end{tabular}
\end{table}

\begin{table}[t]
\centering
\caption{Task-probe rank-1 accuracy for true features vs.\ the no-bridge reference, at both budgets. The small model's IMU probe collapses to chance at the $10\times$ budget; the 230M model's strengthens.}
\label{tab:task}
\begin{tabular}{llccccc}
\toprule
 & & & \multicolumn{2}{c}{33.6M} & \multicolumn{2}{c}{230M} \\
Probe & Budget & Chance & true & none & true & none \\
\midrule
IMU & standard & 0.17 & 0.80 & 0.17 & 0.72 & 0.16 \\
IMU & 10$\times$ & 0.17 & 0.17 & 0.17 & 0.89 & 0.16 \\
AUDIO & standard & 0.02 & 0.33 & 0.03 & 0.31 & 0.02 \\
AUDIO & 10$\times$ & 0.02 & 0.36 & 0.03 & 0.61 & 0.02 \\
\bottomrule
\end{tabular}
\end{table}

\section{What softens}
\paragraph{Prefix tuning flips sign.}
On the small model, feature-conditioned prefix tuning strictly degraded BPB (-0.0048); at 230M it improves it (+0.0614) --- yet remains the weakest adapter by a wide margin against the $\sim$+0.0856 weight-space rows. The original claim is therefore correctly scoped to small models trained without prefixes, while the paper's ordering --- weight-space conditioning dominates token-space conditioning --- holds at both scales.

\paragraph{The diversity--BPB tradeoff softens.}
On the small model, enforcing input-dependent weights at $\lambda=0.20$ erased the BPB gain entirely; at 230M the same penalty retains 89\% of it (+0.0767 of +0.0860). The non-monotonicity the paper reports at the largest tested weight (cross-input cosine reversing at $\lambda=0.20$) appears at the standard budget at 230M as well (0.60 at $\lambda=0.10$ vs.\ 0.72 at $\lambda=0.20$), but at the $10\times$ budget the reversal disappears and cross-input cosine falls to 0.36 at the largest weight (Table~\ref{tab:div}). The penalty required for input-dependence grows with both budget and scale: on the small model $\lambda=0.01$ loses its effect at the long budget, while at 230M penalties below $\lambda=0.10$ never produce input-dependent weights at either budget.

\begin{table}[t]
\centering
\caption{Cross-input cosine of generated weights (IMU; lower = more input-dependent) across the diversity sweep, at both scales and budgets.}
\label{tab:div}
\begin{tabular}{rcccc}
\toprule
$\lambda$ & 33.6M std & 33.6M $10\times$ & 230M std & 230M $10\times$ \\
\midrule
0.00 & 0.97 & 0.99 & 0.99 & 0.98 \\
0.01 & 0.92 & 0.99 & 0.98 & 0.98 \\
0.05 & 0.65 & 0.31 & 0.96 & 0.98 \\
0.10 & 0.44 & 0.47 & 0.60 & 0.68 \\
0.20 & 0.55 & 0.17 & 0.72 & 0.36 \\
\bottomrule
\end{tabular}
\end{table}

\section{What inverts}
The sharpest warning in the original paper was budget-dependent: at $10\times$ training, the small model's unregularized IMU bridge became behaviorally indistinguishable from no bridge at all (rank-1 0.17, equal to chance), having scored 0.80 at the standard budget. At 230M the same protocol \emph{strengthens} conditioning: IMU rank-1 rises from 0.72 to 0.89 (top-2 1.00), and audio rises from 0.31 to 0.61 against a 0.02 chance level, with every control still at chance. Prolonged task-aligned training is an ally of conditioning at this scale, not a failure mode. The weight-space collapse under pure text loss still occurs at 230M (Table~\ref{tab:div}, $\lambda=0$); what changes is that label-paired training at scale maintains input dependence without regularization over the budgets tested.

\section{Implications}
Three practical updates for anyone building on the paper. First, the evaluation discipline (capacity-matched controls, task-aligned probes) transfers verbatim and remains necessary: BPB is exactly as blind at 230M as at 33.6M. Second, the mechanism's failure modes shift in the favorable direction from the 33.6M model to the 230M one --- the task-conditioning collapse threshold moves outward and the cost of enforcing input-dependent weights shrinks. Two caveats bound that statement: the comparison confounds scale with pretraining (the small model was trained from scratch, LFM~2.5 is pretrained), and it concerns task-conditioned training only --- collapse under pure text loss is unchanged at 230M. Even so, claims about weight-space conditioning's fragility measured at one scale should be read as lower bounds. Third, token-space and weight-space conditioning both improve with scale, but the gap between them persists; weight-space modulation remains the stronger interface across every configuration tested.

\section*{Reproducibility}
The port is additive: the published harness reproduces the original paper bit-identically, and the same runner produces every number here via \texttt{--checkpoint hf:LiquidAI/LFM2.5-230M}. Artifacts: the standard-budget suite (35 experiments) and the $10\times$ suite (33 of 35; the two remaining seed-stability repetitions were not run) ship alongside this note as \texttt{summary-lfm.json} and \texttt{summary-lfm-scaling.json} in the code repository, with the original paper's artifacts unchanged. Code: \url{https://github.com/kortexa-ai/legolm.paper}.

\begin{thebibliography}{9}
\bibitem{penov2026bridges} F.~Penov. Conditional LoRA Bridges for Modular Sensor Adaptation of Frozen Small Language Models. Preprint, 2026. \url{https://github.com/kortexa-ai/legolm.paper}.
\bibitem{liquidai2025lfm2} Liquid~AI. LFM2.5-230M. Model release, 2025--2026. \url{https://huggingface.co/LiquidAI/LFM2.5-230M}.
\end{thebibliography}

\end{document}
